% Original Author: Chip Mappus (cmappus@gatech.edu)
% Last Rev:

% Example slides for yet another unofficial Georgia Tech beamer theme
% See http://latex-beamer.sourceforge.net/ for more information on
% beamer.
%Use the following for text over image animation
%http://tex.stackexchange.com/questions/119788/overlay-multiline-text-in-the-image-in-beamer
%or use Ipe tool to make such images with text over them
% Modified by Abhinav Narain 
\documentclass{beamer}

\mode<presentation>{\usetheme{Atlanta}}
\usepackage{natbib}
\usepackage{graphicx}
%\usepackage[footnotesize]{subfigure}
\usepackage[latin1]{inputenc}
% uncomment for Times font
%\usepackage{times}
%\usepackage[T1]{fontenc}
\usepackage{appendixnumberbeamer}
\usepackage{rotating}
\usepackage{color}
\usepackage{tikz}
\title[] {Characterizing Wireless Pathologies in Home Networks}

\author[] 
{Abhinav Narain, \\
Prof. Nick Feamster,\\
Prof. Alex Snoeren (UCSD)}

\institute[PhD Qualifier Presentation]
{College of Computing\\
School of Computer Science\\
 Georgia Institute of Technology}

\begin{document}

{
\usebackgroundtemplate{\includegraphics[width=\paperwidth]{images/logos/Georgia-Tech-Insignia-Watermark-1200x1100}}
\begin{frame}
  \titlepage 
\end{frame}
}

\section{Introduction}
\begin{frame}
  \frametitle{Why study \textit{Wifi} in homes?}
  \begin{itemize} 
  \item \textit{Wifi} technology is the most used for connectivity to the Internet
  \item Home Networks are \textit{unmanaged} and \textit{unorganized}
  \item Presence of many devices that can potentially cause disruption
    in \textit{Wifi} communication e.g. microwave, cordless phone, zigbee
    devices etc.
  \item How does an average home user know the causes of poor
    performance of \textit{Wifi} at home?
  \end{itemize}
   \begin{overlayarea}{\textwidth}{.4\textheight}
     \centering
     \includegraphics[width=.4\textwidth]{./images/wireless_frust.png}
   \end{overlayarea}
\end{frame}

\begin{frame}
  \frametitle{Problems with unmanaged networks} 
    \begin{itemize}
      \item Nearly unusable network performance in many homes
      \item Connectivity problems due to proximity
      \item Co-channel Interference between nodes
      \item TCP might interact poorly because of persistent packet losses
    \end{itemize}
\end{frame}


\begin{frame}
  \frametitle{Troubleshooting Wireless in Home Networks}
    \begin{itemize} 
    \item Anecdotal evidence suggest home networks are small but unpredictable
    \item Spectrum analyzers for detecting broadband Interference
    \item \textit{NetAnalyzer} $/$ \textit{Speedtests} for throughput and delay measurements for short intervals
    \end{itemize}
\end{frame}

\begin{frame}
  \frametitle{Understanding Home Networks} 
    \begin{itemize} 
      \item There is lack of understanding of how 802.11 performs in home networks
      \item What does normal wireless traffic look like?
        \begin{itemize}
          \item How many average retransmissions?
          \item How stable are bitrates devices transmit at?
        \end{itemize}
        \item How to differentiate good from bad?
        \item What is the ground truth? Baselines for such statistics
    \end{itemize}
\end{frame}

\begin{frame}
  \frametitle{Research Questions?} 
  \begin{itemize} 
  \item What is the cause of performance degradation?
    \begin{itemize} 
    \item Broadband Interference
    \item Hidden Terminals
    \item Contention period
    \end{itemize}
  \item What is the most persistent/prevalent issue?
  \item Operating regimes for 802.11 standard
    \begin{itemize} 
    \item $a/b/g/n$ in the same environment  
    \item Maximum operational throughput  in real deployment
    \end{itemize}
  \end{itemize}
\end{frame}

\begin{frame}
  \frametitle{Theoretical vs Practical Approach} 
  \begin{itemize}
  \item Model
    \begin{itemize}
    \item Model for RF interference, signal attenuation of different devices
    \item Model different kinds of construction material in homes (unlike campus/enterprise network)
    \item Cannot model user behavior 
    \end{itemize}
  \item Perform simulation
    \begin{itemize}
    \item  Get results based on assumptions of the model 
    \end{itemize}
  \end{itemize}
\end{frame}

\begin{frame}
  \frametitle{Theoretical vs Practical Approach} 
    \begin{itemize} 
    \item But reality is far different from what we can infer from these simulations!
    \item It is hard to predict$/$model variable channel conditions because of multipath
    \item Need actual testbed!
    \item We need actual measurements to understand conditions of wireless medium and identify pathologies
    \item Many publications assume such problems and give solution
    \end{itemize}
\end{frame}

\section{Approach to solving problem}
\begin{frame}
  \frametitle{Our Approach} 
    \begin{itemize} 
    \item Measurements from actual deployment
    \item Perform \textit{passive}, \textit{continuous} measurements from Internet gateway
      \begin{itemize}
      \item Does not introduce traffic overhead into network
      \item Scheduling regular tests from/to user devices is not required
      \item Gives finegrained information for anomaly detection
      \end{itemize}
    \item Complete control over Access Point
    \item No access to home users wireless devices
    \end{itemize}
\end{frame}

\begin{frame}
  \frametitle{Rationale behind our approach} 
    \begin{itemize}
    \item Take advantage of a pre-existing deployment of BISMark
    \item Have uniformity in the measurements because of the same networking stack and hardware on all routers
    \item Cooperation from home users not required
    %\item $\dots$
    \end{itemize}
\end{frame}

\begin{frame}
  \frametitle{Mechanism to Detect wireless pathologies}
  \begin{itemize}
  \item Contention in medium
    \begin{itemize}
    \item Measure the time difference data frame is enqueued and
      dequeued in the hardware queues when the queue is empty.
    \end{itemize}
  \item Broadband Interference
    \begin{itemize}
       \item Measure the frequency of OFDM timing errors at Physical
         layers 
    \end{itemize}
  \item Hidden Terminal
    \begin{itemize}
    \item Measure the number of RTS-CTS frames in the wireless medium
    \end{itemize}
  \end{itemize}
\end{frame}


\begin{frame}
  \frametitle{Can We Do More?}
    \begin{itemize} 
    \item Constrained to get access to and instrument other home devices 
      \begin{itemize}
      \item Cannot have more infrastructure in one home and get global view of the network
      \end{itemize}
    \item Constrained by Wireless chipset (hardware)
      \begin{itemize} 
      \item \textit{Ath9k} driver does not provide access to fine grained information of \textit{802.11 Contention Window}, \textit{AIFS} (Arbitration Inter-frame Spacing), \textit{TXOP}
      \end{itemize}
    \end{itemize}
\end{frame}

\section{Experiment Design}
\begin{frame}
  \frametitle{Our Peephole: Home Gateway}
  \begin{columns} [onlytextwidth]%[T]
    \begin{column}{.6\textwidth}
      \begin{itemize} 
      \item Hardware
        \begin{itemize} 
        \item Netgear routers
        \item Atheros wireless cards (b/g/n),(a/n)
        \item Dual radio wireless cards
        \end{itemize}
      \item Software
        \begin{itemize} 
        \item OpenWrt firmware
          \begin{itemize} 
          \item white russian and backfire
          \end{itemize}
        \item Modified ath9k driver and kernel subsystem
        \item Custom written userland wireless monitoring tool: \textit{Oculus}
        \end{itemize}
      \end{itemize}
  \end{column}
  \begin{column}{.4\textwidth}
    \includegraphics[width=.8\textwidth]{./images/router.png}
  \end{column}
\end{columns}
\end{frame}

\begin{frame}
  \frametitle{Oculus: What are we Collecting?} 
  \begin{itemize} 
    \item Prepend radiotap headers in a copy of every frame to collect
    \begin{itemize}
       \item Physical Layer, MAC Layer information
       \item Transmit and received timestamps
    \end{itemize}
  \item Collect all data frames
  \item Subset of control and management frames
  \item Measurements data from 8 homes for over 3 months
  \end{itemize}
  \centering
  \includegraphics[width=.8\textwidth]{./images/jig_header.png}
\end{frame}

\section{Challenges}
\begin{frame}
  \frametitle{Research Challenges}
    \begin{itemize}
    \item Limited visibility into the home network
    \item Identification of wireless pathologies using single monitor
      establishing the Ground truth 
    \item Characterizing Wireless conditions$/$behavior
    \item Building Practical and scalable measurement system
      \item Establishing Ground Truth
      \begin{itemize}
        \item Hard to recreate same wireless channel conditions in controlled lab environment
        \item Unlabelled devices and user behavior
      \end{itemize}
    \item Minimal resource consumption on router
      \begin{itemize}
        \item Measurement overhead should not effect own measurements
      \end{itemize}
    \end{itemize}
\end{frame}

\begin{frame}
  \frametitle{Engineering Challenges}
    \begin{itemize} 
    \item Reverse-engineering
      \begin{itemize} 
      \item Interpreting hardware descriptors from ath9k chipset (without programming manual or data sheets)
      \end{itemize}
    \item Instrumenting wireless driver, 802.11 subsystem
      \begin{itemize} 
      \item track queue sizes
      \item transmission timestamps
      \end{itemize}
    \item Custom userland tool to perform light-weight continuous measurements
    \item Remote upgrade might break routers losing deployments
    \item Dealing with buggy Firmware
    \end{itemize}
\end{frame}


\section{Related Work}
\begin{frame}
  \frametitle{Related Studies} 
    \begin{itemize}
    \item Measurement study of 802.11a/b/g networks
    \item WLAN Enterprise Networks
      \begin{itemize}
      \item Sufficient resources
      \item Network Operators survey and design the wireless infrastructure
      \item Jigsaw and Airshark 
      \end{itemize}
    \item Wireless Mesh Network
      \begin{itemize}
      \item \textit{Roofnet} multihop network
      \item \textit{WifiProfiler}: troubleshooting tool for DNS etc services
      \end{itemize}
    \item Home Networks
      \begin{itemize}
      \item  Controlled experiments with multiple monitors in 3 homes
      \end{itemize}
    \item Our study: on 2.4 GHz and 5 GHz; include 802.11 n; has more deployment 
    \end{itemize}
\end{frame}

\section{Observations}
\begin{frame}
  \frametitle{Observations}
    \begin{itemize} 
    \item Measurements data from 20 homes for 2 days    
    \item Characterizing home deployment
       \begin{itemize}
         \item Dense deployment (Apartment type)
         \item Sparse deployment (Single family homes)
       \end{itemize}
    \item Density of devices in homes
     \item 2.4 GHz : Maximum of over 350 wireless devices and 60 Access Points
     \item 5.8 GHz : Maximum of 20 wireless devices and 3 Access Points
    \end{itemize}
\end{frame}
\section{Per Home Observation}
\begin{frame}
  \frametitle{How Poor can the Wireless perform?}
  \begin{tikzpicture}
    \node[anchor=south west,inner sep=0] (image) 
    at (0.5,0)
         {\includegraphics[width=.8\textwidth]{./images/bitrates.png}};
         \begin{scope}[x={(image.south east)},y={(image.north west)}]
           \node<2-> at  (0.78,0.75) {%
             \begin{itemize}
                \item \textcolor{red}{55\% unicast data frame transmissions are at 1 Mbps!}
             \end{itemize}
           };
         \end{scope}
  \end{tikzpicture}
\end{frame}

\begin{frame}
  \frametitle{Asymmetry in Wireless Channel}
  \begin{tikzpicture}
    \node[anchor=south west,inner sep=0] (image) 
    at (0.3,0)
         {\includegraphics[width=.9\textwidth]{./images/asymm.png}};
         \begin{scope}[x={(image.south east)},y={(image.north west)}]
           \node<2-> at  (0.85,0.53) {%
           \begin{overlayarea}{.4\textwidth}{.2\textheight}
             \only<1> \ \textcolor{red}{Asymmetry}
             \end{overlayarea}
           };
         \end{scope}
  \end{tikzpicture}
\end{frame}

\section{Observations Across homes }
\begin{frame}
  \frametitle{Inferring Hidden Terminals in home Wireless Network}
  \begin{tikzpicture}
    \node[anchor=south west,inner sep=0] (image) 
    at (0.3,0)
         {\includegraphics[width=.8\textwidth]{./images/tty1.png}};
         \begin{scope}[x={(image.south east)},y={(image.north west)}]
           \node<2-> at  (0.78,0.75) {%
             \begin{overlayarea}{.4\textwidth}{.2\textheight}
               \only<1> \ \textcolor{red}{\textit{HT-Minstrel} rate algorithm  uses RTS-CTS  following every unsuccessful transmission!}
             \end{overlayarea}
           };
         \end{scope}
  \end{tikzpicture}
\end{frame}


\begin{frame}
  \frametitle{Background on Calculation of Contention Delay} 
    \begin{itemize} 
    \item Total Transmission time $=$ Queueing Delay $+$ Channel Access Delay 
    \item Empty Queue$ => $Queueing Delay $ =0 $
    \item Total Transmission time (Contention Delay) $=$Channel Access Delay 
    \end{itemize}
    \begin{itemize}
    \item \textit{802.11 e } Primer
    \item Four Access Classes: Video, Voice, Best Effort, Background Traffic
    \item Four hardware queues in the Wireless NIC 
    \end{itemize}
\end{frame}
\begin{frame}
  \frametitle{Variation of Contention Delay(B.E) with Neigborhood Access Points}
  \begin{tikzpicture}
    \node[anchor=south west,inner sep=0] (image)
    at (0.0,0)
         {\includegraphics[width=.8\textwidth]{./images/contention_delay.png}};
         \begin{scope}[x={(image.south east)},y={(image.north west)}]
           \node<2-> at (0.58,0.75) {%
           \begin{overlayarea}{.4\textwidth}{.2\textheight}
             \only<1> \ \textcolor{red}{Contention delays are uncorrelated with number of access points}
             \end{overlayarea}
             % \begin{overlayarea}{\textwidth}{4cm}
             %    \only<1> $55\%$ transmissions are at 1 Mbps 
             % \end{overlayarea}     
           };
         \end{scope}
  \end{tikzpicture}
\end{frame}

\begin{frame}
  \frametitle{Percentage Retransmissions with Number of Devices }
  \begin{tikzpicture}
    \node[anchor=south west,inner sep=0] (image)
    at (0.0,0)
         {\includegraphics[width=.7\textwidth]{./images/var_retx.png}};
         \begin{scope}[x={(image.south east)},y={(image.north west)}]
           \node<2-> at (0.58,0.75) {%
           \begin{overlayarea}{.4\textwidth}{.2\textheight}
             \only<1> \ \textcolor{red}{Retransmissions are independent of number of devices}
             \end{overlayarea}
           };
         \end{scope}
  \end{tikzpicture}
\end{frame}

\section{Controlled Experiment}
\begin{frame}
  \frametitle{Detection of Broadband Interference}
  \begin{tikzpicture}
    \node[anchor=south west,inner sep=0] (image)
    at (0.0,0)
         {\includegraphics[width=.9\textwidth]{./images/mwave.png}};
         \begin{scope}[x={(image.south east)},y={(image.north west)}]
           \node<2-> at  (0.58,0.85) {%
             \begin{itemize}
             \item \textcolor{red}{Microwave increases the Error decoding 802.11
               frames }
             \end{itemize}
           };
         \end{scope}
  \end{tikzpicture}
\end{frame}


\section{Conclusion}
\begin{frame}
  \frametitle{Conclusion and Future Work} 
    \begin{itemize} 
    \item Designed, Implemented measurement system for home networks        
    \item Analysis of collected data of BISMark deployment
    \item Deviced mechanism to measure broadband Interference using
    commodity hardware
    \item Collect transport layer traces to correlate poor performance
    across layers
    \item Design in-kernel aggregiate data structures for accurate broadband
    Interfering source identification
    \end{itemize}
\end{frame}

\begin{frame}
  \frametitle{Publication/Poster/Talks} 
    \begin{itemize} 
    \item Talk at IS4CWN (International Summit for Community Wireless
      Networks), Berlin. Oct, 2013
    \end{itemize}
\end{frame}

\end{document}

\begin{frame}
  \frametitle{}
  \begin{tikzpicture}
    \node[anchor=south west,inner sep=0] (image) 
    at (0.5,0)
         {\includegraphics[width=.5\textwidth]{./images/}};
         \begin{scope}[x={(image.south east)},y={(image.north west)}]
           \node<2-> at  (0.78,0.75) {%
             \begin{itemize}
             \item 
             \end{itemize}
             % \begin{overlayarea}{\textwidth}{4cm}
             %    \only<1> $55\%$ transmissions are at 1 Mbps 
             % \end{overlayarea}     
           };
         \end{scope}
  \end{tikzpicture}

\end{frame}

\begin{frame}
  \frametitle{} 
    \begin{itemize} 
    \end{itemize}
\end{frame}

% Local Variables:
% mode: latex
% mode: flyspell
% mode: auto-fill
% TeX-master: t
% End:
%   \begin{block}{}
%     \begin{enumerate}
%     \item 
%     \end{enumerate}
%   \end{block}

%\begin{frame}
%  \frametitle{}
%    \begin{itemize} 
%    \end{itemize}
%\end{frame}
